\documentclass[11pt]{article}
\usepackage[review]{acl}
\usepackage{times}
\usepackage{latexsym}
\usepackage{amsmath,amssymb}
\usepackage{booktabs}
\usepackage{microtype}
\usepackage{array}
\usepackage{tabularx}
\usepackage{url}
\urlstyle{tt}
\usepackage[utf8]{inputenc}
\title{Program-Conditioned Relation Tests for Numerical Error Ranking: Evidence and Boundaries}
\author{Anonymous NAACL submission}
\begin{document}
\maketitle
\begin{abstract}
Can controlled perturbations rank numerical errors? We study relation-only PECR on FinQA. Construction is program-conditioned: gold executable programs and transformed program results determine expected signs. Scoring is numeric-answer-blind: it uses frozen signs and parsed model answers, not numeric gold answers or correctness labels. On 139/157 selected TEST items (94 errors), relation risk reaches error AUROC 0.7337 (95\% CI [0.6615, 0.7991]) and AUPRC 0.8164 (error prevalence 0.676). It exceeds confidence and both Any-change orientations, but its AUROC gain over the fixed-majority-direction baseline is +0.0500 (95\% CI [-0.0092, +0.1215]), not a clear advantage. Eighteen unscorable items were later all labeled incorrect. Two post-hoc submitted construction forms agree on 14/20 INVALID, 2 UNCERTAIN, and 4 VALID judgments; answer judgments agree on 28/28 unique items. The project lead confirms that the second form was independently completed by a lab classmate; the reviewer-ID field was entered incorrectly. Because this identity history is project-attested rather than externally authenticated, agreement is reported descriptively. In a selected leave-out sensitivity excluding the 14 jointly flagged constructions, the relation-versus-majority AUROC gap falls to +0.0087 (95% CI [−0.0187, +0.0485]). These reviews do not establish construct validity. The TEST label rule was frozen after outputs and before label generation, but ordering is not independently verified. This polarity-corrected analysis is not an unchanged preregistered confirmation. CST provides no independent corroboration. We report a bounded finding, not end-to-end oracle-free construction, deployment readiness, or a shared PECR--CST mechanism.
\end{abstract}

\section{Introduction}
Confidence and answer agreement are useful but imperfect indicators of whether a numerical answer is correct \citep{guo-etal-2017-calibration,wang-etal-2022-self-consistency,kuhn-etal-2023-semantic-uncertainty}. A model may be confident and wrong, or may change its answer under perturbation for reasons unrelated to whether the change is appropriate. We ask whether consistency with known relations across controlled changes can rank errors in an original answer.

We evaluate Program-Executable Counterfactual Reliability (PECR) on FinQA, a benchmark for numerical reasoning over financial text \citep{chen-etal-2021-finqa}. We focus on a relation-only score and fixed cohort. The result is positive but bounded: error-positive AUROC is 0.7337 on 139 scorable selected TEST items, exceeding confidence risk and both orientations of an Any-change comparison. Against a fixed majority-direction baseline, however, the AUROC increment is +0.0500 with a paired interval crossing zero.

The information boundary is essential. Construction uses gold programs and program-derived transformed results to validate counterfactual worlds and freeze expected signs. Scoring is blind to numeric gold answers, correctness labels, transformed gold answers, and teacher CEF. We therefore describe PECR as \emph{program-conditioned, numeric-answer-blind scoring}; it is not end-to-end oracle-free. Correctness labels enter only at evaluation.

Our contributions are: (1) an empirical error-ranking result for relation risk on a fixed, partially scorable FinQA TEST subset; (2) an explicit account of oracle use, protocol changes, polarity correction, exclusions, and label provenance; and (3) a construct-validity boundary for CST, not independent confirmation of PECR.

\section{PECR and its information boundary}
\paragraph{Construction.} For each selected question, PECR uses the gold executable program and an operand-to-source mapping to construct numerical counterfactuals. The program is run on transformed operands to verify outcomes and freeze expected signs. Construction is therefore program-conditioned and depends on benchmark-provided executable information.
\paragraph{Prompting.} The model receives the original and transformed question/evidence text for fixed worlds under a frozen prompt. Prompts do not expose gold numerical answers or correctness labels; transformed program outputs are not answer hints.
\paragraph{Scoring.} Responses are parsed under the frozen v2.2 numeric parser. Relation risk compares observed answer movements with expected signs fixed during construction. The scoring code does not read the gold program, gold answer, correctness label, transformed gold answer, or teacher CEF. This no-oracle statement applies to scoring, not construction.
\paragraph{Evaluation.} After score and coverage freezing, original-answer labels are derived by comparing the parser-valid original free-text answer to FinQA TEST `qa.answer` under a separately versioned rule migrated from TRAIN. This is not the official FinQA program-execution metric.

\begin{table}[t]
\centering\small
\begin{tabularx}{\columnwidth}{@{}p{0.19\columnwidth}X p{0.27\columnwidth}@{}}
\toprule
Stage & Information used & Boundary \\
\midrule
Construction & Gold program, operand/source mapping, transformed program results & Program-conditioned; not oracle-free \\
Prompt & Fixed original/transformed question text & No gold numeric answer or correctness label \\
Scoring & Parsed model answers and frozen expected signs & Numeric-answer-blind; no gold program, gold answer, transformed gold, label, or teacher CEF \\
Training & None for relation-only risk & No supervised TEST fitting \\
Evaluation & Frozen labels derived from `qa.answer` & Post-output, pre-label policy; not initial preregistration \\
\bottomrule
\end{tabularx}
\caption{Information boundaries for relation-only PECR. Gold-answer blindness in scoring must not be generalized to construction.}
\label{tab:boundary}
\end{table}

\section{Experimental setup}
\paragraph{Cohort and calls.} The fixed manifest contains 157 FinQA TEST questions and a budget of 471 world requests (three per question). The server deployment was Qwen3.5-4B. Parser-invalid outputs were recorded without retry; transport, model-identity, and global integrity failures remained stop conditions. Slot 40 triggered the earlier parser stop. Before slots 41--471, the continuation rule was versioned so parser-invalid responses were retained as invalid and the fixed sequence continued without replacement. This is a run-time protocol amendment.
\paragraph{Scorable set.} The label-free finalizer required all three worlds and exact normalized unit compatibility, or all units absent; no conversion was applied. This yielded 139/157 (88.5\%) scorable items. Thirteen were excluded for unit-incompatible triples and five for incomplete triples. All 18 were later labeled incorrect, so the scorable subset is outcome-selected and cannot be assumed representative.
\paragraph{Outcome and statistics.} The positive class is \texttt{original\_incorrect}; larger scores indicate more risk. We report AUROC and AUPRC with 95\% percentile intervals from 5,000 source-group bootstrap replicates (fixed seed 20260923). Paired score differences use shared resamples. Error prevalence, 94/139=0.6763, is the no-skill AUPRC reference.

\section{Results}
\begin{table*}[t]
\centering\small
\begin{tabularx}{\textwidth}{@{}Xcc@{}}
\toprule
Risk score (larger = higher error risk) & Error AUROC [95\% CI] & Error AUPRC [95\% CI] \\
\midrule
Relation risk & \textbf{0.7337 [0.6615, 0.7991]} & \textbf{0.8164 [0.7408, 0.8804]} \\
Original-confidence risk ($1-$confidence) & 0.5039 [0.4245, 0.5854] & 0.6873 [0.6039, 0.7685] \\
Literal Any-change fraction & 0.3918 [0.3370, 0.4486] & 0.6291 [0.5407, 0.7154] \\
Any-change complement ($1-$fraction; v1 analyzer direction) & 0.6082 [0.5514, 0.6630] & 0.7351 [0.6532, 0.8099] \\
Fixed majority directions ($-1,+1$; post-hoc control) & 0.6837 [CI not estimated] & 0.7707 [CI not estimated] \\
\bottomrule
\end{tabularx}
\caption{Error ranking on the fixed 139-item scorable TEST subset (94 incorrect, 45 correct). Values come from versioned error-positive analyses; intervals are not estimated for the post-hoc fixed-majority row. The initial analysis used the wrong positive-class polarity for risk scores.}
\label{tab:primary}
\end{table*}

The paired AUROC increment for relation risk is +0.2298 [+0.1328, +0.3289] over confidence risk, +0.3418 [+0.2250, +0.4514] over literal Any-change, and +0.1255 [+0.0740, +0.1782] over its complement. A separately computed post-hoc fixed-majority-direction control ($-1$ for the minus world and $+1$ for the plus world) obtains AUROC 0.6837 and AUPRC 0.7707. Relation risk exceeds it by +0.0500 AUROC, but the paired source-group bootstrap 95\% CI [-0.0092, +0.1215] crosses zero; thus we do not establish an advantage over this stronger direction prior. The original preregistration does not designate this later control as primary. Paired AUPRC increments are +0.1291 [+0.0695, +0.1970], +0.1873 [+0.1108, +0.2666], and +0.0813 [+0.0483, +0.1211], respectively. The preregistration defines the literal change fraction but does not explicitly designate its complement as the primary risk comparator; the original analyzer used the complement. Both directions are reported without post-result selection. The numerical G2 criterion against the original analyzer comparator is met only in this versioned error-positive analysis, not by an unchanged initial implementation; this does not establish superiority over the later fixed-majority-direction control.

\paragraph{Post-hoc construction-flagged sensitivity.} To assess sensitivity without changing the primary set, we recomputed the same metrics after excluding (a) the five AMT/BLK cases singled out by the LLM audit and (b) all 20 purposively selected construction-diagnostic items. These are exploratory leave-out analyses, not judgments that the items are invalid. Excluding the five gives $n=134$, relation AUROC 0.7448, fixed-majority AUROC 0.7016, and paired difference +0.0432 [-0.0111, +0.1076]. Excluding all 20 gives $n=119$ and identical AUROCs of 0.7295 for relation and fixed-majority scores. Importantly, the 20 diagnostics are exactly the 20 items whose expected directions differ from the fixed $(-1,+1)$ prior; the remaining 119 use the majority directions. Therefore, this leave-20-out equality is mechanically expected and is not independent validation or evidence of subgroup performance. The primary 139-item analysis remains unchanged.


\paragraph{Human-review-defined construction sensitivity.} Using the two submitted forms' shared VALID/INVALID/UNCERTAIN classifications, we recomputed the relation and fixed-majority scores after leaving out reviewed items. Excluding the 14 items both forms mark INVALID leaves $n=125$ and reduces the relation-minus-majority AUROC difference to +0.0087 [-0.0187, +0.0485]. Excluding the 14 INVALID plus two UNCERTAIN items leaves $n=123$ and a difference of +0.0092 [-0.0193, +0.0492]. Excluding all 20 purposively selected direction-diagnostic items leaves $n=119$, where the relation and majority-direction AUROCs are identical (0.7295), because the remaining items use the fixed majority directions. These post-hoc, selection-conditioned sensitivities do not revise the primary result or estimate cohort-wide construction validity; they indicate that the apparent incremental signal over the strong direction prior is concentrated in the reviewed diagnostic cases and remains uncertain.

\begin{table}[t]
\centering\small
\begin{tabularx}{\columnwidth}{@{}Xrcc@{}}
\toprule
Post-hoc subset & $n$ & Relation AUROC & Relation $-$ majority [95\% CI] \
\midrule
Primary scorable cohort & 139 & 0.7337 & +0.0500 [-0.0092, +0.1215] \
Exclude 14 reviewer-marked INVALID & 125 & 0.7204 & +0.0087 [-0.0187, +0.0485] \
Exclude 14 INVALID + 2 UNCERTAIN & 123 & 0.7322 & +0.0092 [-0.0193, +0.0492] \
Exclude all 20 purposive diagnostics & 119 & 0.7295 & +0.0000 [0.0000, 0.0000] \
\bottomrule
\end{tabularx}
\caption{Post-hoc construction-review sensitivity. Review-defined exclusions are not population-validity estimates; the diagnostic items were selected because their signs differ from the fixed majority prior.}
\label{tab:human-review-sensitivity}
\end{table}

\paragraph{Separate post-hoc student control.} In a historical ConvFinQA setting, direct GBDT AUROC was 0.8183 [0.7211, 0.9001] and imputed-student AUROC was 0.7474 [0.6449, 0.8408]. The post-hoc direct-minus-imputed difference was +0.0709 [-0.0239, +0.1682], crossing zero. This is not a same-cohort FinQA baseline and does not support imputation superiority.

\section{CST as a construct boundary}
CST is not presented as an equally mature success or shared-mechanism confirmation. In the natural-reverse audit, reverse/original prompt pairs were byte-identical and paired-answer behavior reproduced most apparent reverse-axis signal; natural-reverse is therefore not an independent evidence-direction test. In the strict independent-counter-evidence pilot, the flip contrast was +0.2118 [0.1160, 0.3097], but the placebo flip rate was 0.5141 and the placebo criterion failed; there were only three high-consensus wrong items and independent-score AUROC was 0.624 [0.286, 0.856]. This supports at most a behavioral contrast, not reliable independent ranking. CST full cohort remains deferred.

\paragraph{Post-hoc LLM-assisted audit.} After the primary analysis, we used two configured LLM reviewers to audit (i) a purposively selected 20-item construction diagnostic set and (ii) a label-review queue containing 35 records for 28 unique questions. The construction reviewers saw the question, three constructed worlds, and program-derived construction basis, but not frozen expected signs, evaluated-model answers, risk scores, or correctness labels. Thus this audit was itself program-conditioned. The answer reviewers saw the question, original answer, and gold answer, but not PECR risk scores or frozen correctness labels. This was not independent human adjudication: one reviewer used the same Qwen3.5-4B model family as the evaluated system, while the other was routed as gpt-6-luna through Codex; the underlying provider and model identity for that route are not independently established by the retained records. The audit is post-hoc, and its final comparison rules were corrected after V2 outputs had been inspected; it is exploratory rather than blind or preregistered.

On construction, the configured gpt-6-luna reviewer supplied 18 determinate judgments per world among 20 items, all matching frozen signs, while marking two unresolved. This match is not a construction-validity rate: its rationales also flagged contradictory evidence. Qwen supplied parseable judgments for 12/20 items; its judgments matched 2/12 minus-world signs and 5/12 plus-world signs. Two Qwen batches were malformed or truncated and were not retried. Among the 12 items with both reviewers' outputs, they agreed on direction for 1 minus-world and 5 plus-world cases. The diagnostic set was purposively selected and cannot estimate error prevalence for all 157 questions. In particular, AMT/2003 retained a stated remaining area of 10,000 despite transformed arithmetic implying 12,000 or 8,000; AMT/2010 and BLK/2012 changed totals while component rows remained inconsistent; and BLK/2015 contained conflicting rate representations. These are unresolved construct concerns, not corrected cases.

For answer labels, the two reviewers gave the same resolved judgment on 22/28 unique questions and disagreed on six. All 22 unanimous judgments matched the frozen labels, so the prespecified-for-this-audit unanimous-only sensitivity procedure changed no labels and left the primary metrics unchanged. This does not establish label truth: the audit was post-hoc, reviewers saw gold answers, six cases were disputed, and agreement between configured LLM reviewers is not human adjudication or independent validation. No frozen label, score, manifest, ledger, or primary-analysis artifact was altered.

\paragraph{Post-hoc submitted review forms and provenance.} Two completed form sets cover the 20 purposively selected construction diagnostics and the answer/gold queue (35 records representing 28 unique questions). The project lead confirms reviewer 2 is a lab classmate who independently completed the forms; the reviewer-ID field in the submitted CSV is erroneous. The original forms are preserved unchanged, and the corrected identity is recorded as project-lead-attested metadata rather than written into the source CSV. The project history contains an earlier instruction phrased as a request to copy and relabel forms; the project lead subsequently clarified that this did not describe how reviewer 2 actually completed the forms. We therefore report item-level categorical agreement between two human reviewers descriptively. Since identity and completion circumstances are not externally authenticated, we do not present these figures as a formal reliability estimate or independent replication. On construction, the forms agree on both independently derived directions for 20/20 cases, overall VALID/INVALID/UNCERTAIN status for 20/20 (4/14/2), financial coherence for 20/20, and prompt-relation consistency for 20/20; source-operand mapping agrees on 17/20, with disagreements at AMT/2003/page\_27.pdf-1, BLK/2015/page\_124.pdf-4, and FIS/2006/page\_31.pdf-2. The judgments are from a purposively selected diagnostic set and do not estimate cohort validity. On answer review, the forms agree on CORRECT/INCORRECT/AMBIGUOUS for 28/28 unique items (22/4/2), while difference-category agreement is 26/28. Ten questions overlap the primary 139; eight have shared determinate judgments, and the selected label sensitivity changes one of those relative to the frozen label. Replacing only those eight shared determinate labels, while retaining the fixed cohort and all ambiguous/unreviewed labels, yields relation AUROC 0.7380 (versus 0.7337 primary) and relation-minus-majority AUROC +0.0484 [-0.0098, +0.1189]. This remains a post-hoc sensitivity, not a corrected primary result or proof of label truth. Record counts and unique-question counts are distinct.
\section{Protocol chronology and analysis provenance}
The raw run records were retained. Slot 40 triggered the earlier parser stop; before slots 41--471, parser-invalid outputs were recorded and the fixed sequence continued without retries or replacements. Thus the complete run did not use one unchanged stop rule.

After model collection, a label-free finalizer froze the 139 scorable IDs, scores, coverage, and same-cohort controls. A separate TEST correctness policy was then frozen and synthetic tests passed before labels were generated. The policy compares parser-valid original free-text answers to `qa.answer`, uses the frozen TRAIN numeric rule, and covers all 157 items. It is post-output and pre-label, not an initial preregistration; it is not called the official FinQA score.

The continuation receipt reports that fixed slots 41--471 completed at 2026-09-24 16:07:02.916Z (431 attempts, no retries or fallback); the final raw-ledger SHA-256 is `789fc8c67eefdd92392c1893af2d796a94cc91efa1cc60da2fd7a7f26ad1cf3f`. The UTC-converted artifact sequence then places the label-free receipt at 16:09:02.929Z; TEST policy at 16:22:40.791Z; synthetic tests at 16:23:40.331Z; pre-label checkpoint at 16:23:40.361Z; label artifact at 16:23:40.688Z; initial analysis output at 16:26:22.951Z; and error-positive V2 output at 16:55:57.181Z. The policy/checkpoint record says TEST gold had not been read and labels had not been generated. Filesystem times and self-reported receipt times are not trusted timestamp attestations; no independent runtime log records the exact first read of `qa.answer` or analysis process start times. Available records support—but cannot independently verify—the freeze-before-label ordering. The policy was frozen after model outputs and before label generation, not as part of the initial preregistration. At audit time, the host reported NTP synchronization; this present-time status does not establish that the clock was continuously synchronized during the historical events. Full artifact SHA-256 values are listed in the versioned UTC chronology audit.

\begin{table*}[t]
\centering\scriptsize
\begin{tabular}{@{}p{0.31\textwidth}p{0.28\textwidth}p{0.34\textwidth}@{}}
\toprule
Event & UTC & Artifact SHA-256 prefix \\
\midrule
Slots 41--471 completion & 2026-09-24 16:07:02.916Z & \texttt{257b0b75e9a2} (receipt) \\
Raw ledger finalized & 2026-09-24 16:07:02.905Z & \texttt{789fc8c67eef} \\
Label-free cohort/score freeze & 2026-09-24 16:09:02.929Z & \texttt{597c05953562} \\
TEST label policy written & 2026-09-24 16:22:40.791Z & \texttt{7196148d4f9e} \\
Synthetic tests / pre-label checkpoint & 2026-09-24 16:23:40.331Z / 16:23:40.361Z & \texttt{aa726bc8d975} / \texttt{dafe03f207df} \\
Label artifact appears & 2026-09-24 16:23:40.688Z & \texttt{899db5256c0c} \\
Initial / error-positive V2 output files & 2026-09-24 16:26:22.951Z / 16:55:57.181Z & \texttt{2e7661de5d20} / \texttt{3400561f2292} \\
\bottomrule
\end{tabular}
\caption{Record-supported chronology; times are filesystem modification times except the run completion time reported by its receipt. Hash prefixes identify artifacts but do not attest to their creation time. The complete hashes and caveats are in the accompanying chronology audit.}
\label{tab:chronology}
\end{table*}

\section{Limitations}
First, the result applies to one Qwen3.5-4B deployment and a selected subset of FinQA TEST. Eighteen of 157 items were unscorable and all 18 were later labeled incorrect. Outcome-associated missingness makes the 139-item result unsuitable as an unbiased estimate for the full manifest; we do not impute a primary full-cohort AUROC. AUPRC 0.8164 should be interpreted against the high error prevalence, 0.6763.

Second, PECR construction is gold-program-conditioned. Although scoring is numeric-answer-blind, counterfactual construction and expected signs depend on the gold executable program and transformed program results. This is not end-to-end oracle-free and does not establish a deployable intervention generator.

Third, the initial implementation had an outcome-polarity error, slot-40 continuation handling was amended, and the TEST free-text correctness policy was frozen after output collection but before label generation. The initial analysis is preserved. Available artifacts support but cannot independently verify the chronology of the pre-label checkpoint relative to first gold access; the first `qa.answer` read has no independently timestamped event record. The result is not an unmodified preregistered confirmation.

Fourth, the post-hoc review forms have a provenance limitation. The project lead confirms reviewer 2 was a lab classmate who independently completed the forms; the submitted reviewer-ID field was erroneous and is not silently altered in the archived source files. The project history contains an earlier copy-and-relabel instruction, later clarified by the project lead as not describing how the second review was actually completed. We therefore report observed agreement between two human reviewers, with identity/completion provenance attributed to the project lead and not externally authenticated; it is not a formal reliability estimate or independent replication. The reviewers agreed on 4/20 constructions as valid, 14 invalid, and 2 uncertain; they marked financial coherence identically for 20/20. The 20 items were purposively selected and exactly comprise cases whose expected signs differ from the fixed-majority prior, so these judgments do not estimate cohort validity. The answer/gold forms cover 35 records representing 28 unique questions; the reviewers agreed on CORRECT/INCORRECT/AMBIGUOUS for all 28 unique questions, while difference-category agreement was 26/28. Ten questions overlap the primary 139, eight have shared determinate judgments, and one differs from the frozen label. The AUROC sensitivity (0.7380 versus 0.7337) is selected and post hoc, not a corrected primary analysis or proof of label truth. Neither form resolves AMT/BLK concerns or provides external construct validation.

Fifth, the historical direct-GBDT comparison is post hoc and from another ConvFinQA setting; its paired interval crosses zero and gives no basis to claim imputation superiority. CST provides no independent corroboration: natural-reverse is mechanically non-independent and the strict pilot fails its placebo gate. No shared mechanism or general reliability capability is established.

\section{Conclusion}
A program-conditioned, numeric-answer-blind relation score ranked original-answer errors on a fixed 139-item scorable FinQA TEST subset, with error-positive AUROC 0.7337 and positive paired differences against confidence and both orientations of Any-change. Superiority over the fixed-majority-direction control is not established (AUROC difference +0.0500, 95\% CI [-0.0092, +0.1215]); moreover, the difference shrinks to +0.0087 [-0.0187, +0.0485] when excluding the 14 purposively reviewed items both forms judged INVALID. Two reviewers independently completed the post-hoc forms according to the project lead; reviewer-identity provenance is user-attested, with an earlier copy-and-relabel instruction retained as a disclosed historical ambiguity. This bounded finding is qualified by outcome-associated exclusions, protocol amendments, polarity correction, and label chronology that is record-supported but not independently verifiable. CST remains a construct boundary rather than independent confirmation. A second-model replication and a genuinely independent CST preflight are possible future studies, not evidence in this paper.

\bibliographystyle{acl_natbib}
\bibliography{references}
\end{document}
